\section{Implementation in this Repository}
\label{sec:impl}

This section grounds the pattern in the repository that produced this
paper.\footnote{\url{https://github.com/noheton/f-ai-r}} The intent is
not to prescribe a stack: any version-control system with a build
target, a structured-data format that supports PROV-O, and a way to
version prompt files would do. The intent is to show that the cost is
modest in practice.

\subsection{Layout}

\begin{description}
    \item[\texttt{paper/}.] \LaTeX{} sources, bibliography, Makefile,
        figure sources. Two manuscripts share one bibliography:
        \texttt{main.tex} (long form) and
        \texttt{main-condensed.tex} (venue length). Sections live under
        \texttt{paper/sections/} and are \texttt{\textbackslash input}'d
        into the main file.
    \item[\texttt{doc/}.] Methodology, FAIR notes,
        human-AI collaboration policy, research protocol, submission
        plan, append-only logbook, the PROV-O graph in Turtle, and a
        Mermaid mirror of the graph for human readers. Vendored sources
        sit under \texttt{doc/sources/} when the licence permits.
    \item[\texttt{agents/}.] One markdown file per role: orchestrator,
        scientific-writer, source-analyzer, fair-aligner,
        provenance-curator, layout-scrutinizer, readability-reviewer,
        illustration, condenser, research-protocol. Each prompt has the
        same five sections (Role; Primary-artifact consistency;
        You do; You do not; Output) so a different harness can ingest
        it with minimal adaptation.
    \item[\texttt{scripts/}.] A small Python module
        (\texttt{build\_provenance\_site.py}) that parses
        \texttt{doc/provenance.ttl} via \texttt{rdflib} and renders it
        into a static site, and a \texttt{requirements.txt} for the CI.
    \item[\texttt{.github/workflows/}.]
        \texttt{build-paper.yml} compiles both manuscripts on every
        push and uploads the PDFs as artefacts.
        \texttt{pages.yml} validates the Turtle, builds the site, and
        (on \texttt{main}) deploys via the official GitHub Pages
        actions.
    \item[Root.] \texttt{CITATION.cff}, \texttt{codemeta.json},
        \texttt{.zenodo.json}, dual licensing
        (\texttt{LICENSE} for code, \texttt{LICENSE-paper} for prose),
        \texttt{CLAUDE.md} for the harness rules.
\end{description}

\subsection{Provenance schema}

The schema is in \S\ref{sec:bg-prov} and the canonical file is
\texttt{doc/provenance.ttl}. The seed graph at the time of writing
contains, in addition to the \texttt{fair2r:} class declarations, one
\texttt{prov:Activity} per pipeline stage that has run, one
\texttt{fair2r:Source} per BibTeX entry, one \texttt{fair2r:Prompt}
per file under \texttt{agents/}, and one \texttt{fair2r:Claim} per
named claim in this manuscript. Appendix~\ref{app:fair-cells} shows
the corresponding F(AI)\textsuperscript{2}R cell mapping; the seed
graph at the time of submission contains 275 triples.

\subsection{Continuous integration}

\paragraph{Paper build.} The \texttt{build-paper.yml} workflow runs on
every push that touches \texttt{paper/} and on every pull request
likewise. It uses the \texttt{xu-cheng/latex-action@v3} container,
which ships a maintained TeX Live distribution; it compiles both
\texttt{main.tex} and \texttt{main-condensed.tex} with
\texttt{latexmk -pdf -interaction=nonstopmode -halt-on-error} and
uploads both PDFs as branch-tagged artefacts. The workflow is
deliberately silent: build failures are visible in the Actions tab and
not auto-issue-filed, because a build failure under this discipline is
typically a sign that the writer and the curator have desynchronised.

\paragraph{Site build and deploy.} The \texttt{pages.yml} workflow
runs on pushes to \texttt{main} and to working branches.
\textsc{Validate} stage parses \texttt{doc/provenance.ttl} via
\texttt{rdflib} and fails on any malformed Turtle.
\textsc{Build} stage runs the static-site generator.
\textsc{Deploy} stage runs only on \texttt{main} and uses the official
\texttt{actions/configure-pages}, \texttt{actions/upload-pages-artifact},
\texttt{actions/deploy-pages} chain. Working branches see the build
output as a downloadable artefact, never as a public deploy.

A consent gate analogous to the one in
\texttt{noheton/Obscurity-Is-Dead} is not yet in place because the
present manuscript is the consent record; we revisit this in
\S\ref{sec:discussion}.

\subsection{The clone-and-replay loop}

A reader who clones the repository, installs a TeX Live distribution
and the dependencies in \texttt{scripts/requirements.txt}, runs
\texttt{make -C paper pdf}, runs \texttt{python scripts/build\_provenance\_site.py},
and opens \texttt{\_site/index.html}, has reproduced the paper, the
provenance graph, and the public site. The audit pass for
\textbf{Reusable} is, in our setup, that loop run on a clean machine
and compared against the published artefact. Any difference is a
defect.
